\newcommand{\SchoolWorksheetSubject}{Математика}
\newcommand{\SchoolWorksheetGrade}{5 класс}
\newcommand{\SchoolWorksheetSection}{Натуральные числа}

\input{../../../../common/Latex/school-worksheet-common.sty}
\ifdefined\workbookbw
  \input{../../../../common/Latex/school-worksheet-bw.sty}
\else
  \input{../../../../common/Latex/school-worksheet-color.sty}
\fi

\usepackage{pdfpages}

\hypersetup{
  pdftitle={Натуральные числа. Рабочая тетрадь ученика. 5 класс},
  pdfsubject={Основной маршрут M01--M07},
  pdfauthor={Школьные учебные материалы}
}

\renewcommand{\contentsname}{Содержание}
\setcounter{tocdepth}{1}

\newcommand{\technicalblankpage}{%
  \clearpage
  \thispagestyle{empty}%
  \null
  \clearpage
}

\newcommand{\includemodule}[3]{%
  \includepdf[
    pages=-,
    fitpaper=true,
    pagecommand={},
    addtotoc={1,section,1,{#2},#1}
  ]{../../output/pdf/#3-student-\workbookmode.pdf}%
}

\begin{document}
\mainfontsize

\thispagestyle{empty}
\begin{center}
  \vspace*{24mm}

  {\sffamily\bfseries\fontsize{31pt}{37pt}\selectfont
   \color{Primary}Натуральные числа\par}

  \vspace{5mm}

  {\sffamily\fontsize{19pt}{24pt}\selectfont
   \color{Secondary}Рабочая тетрадь ученика\par}

  \vspace{3mm}

  {\sffamily\Large 5 класс\par}

  \vspace{22mm}

  \begin{tcolorbox}[
    width=0.78\textwidth,
    colback=PrimaryLight,
    colframe=Primary,
    boxrule=0.9pt,
    arc=2mm,
    left=7mm,right=7mm,top=6mm,bottom=6mm
  ]
    \centering\sffamily
    Семь учебных модулей: от цифр и разрядов\\
    до полного упорядоченного перебора
  \end{tcolorbox}

  \vfill

  \begin{tabularx}{0.86\textwidth}{@{}X@{}}
    \textbf{Фамилия, имя:}\ \hrulefill\\[9mm]
    \textbf{Класс:}\ \answerblank[42mm]
    \hfill
    \textbf{Учебный год:}\ \answerblank[45mm]
  \end{tabularx}

  \vspace{15mm}
\end{center}

\clearpage
\thispagestyle{firstpage}
\tableofcontents

\vspace{1.2em}

\begin{notebox}[Как работать со сборником]
  \begin{enumerate}
    \item Выполняйте модули по порядку и записывайте решения прямо в листах.
    \item В конце каждого модуля пройдите контрольную точку и проверьте себя.
    \item Дальнейший маршрут определите вместе с учителем: следующий модуль,
      дополнительная практика, коррекция или усложнение.
  \end{enumerate}
\end{notebox}

\begin{reflectionbox}[Условные обозначения]
  \textbf{M01--M07}~--- основные учебные модули.
  Дополнительные карточки учитель выдаёт отдельно по результатам контрольной
  точки.
\end{reflectionbox}

\includemodule{module-m01}{M01. Цифры, числа и разряды}{M01}
\includemodule{module-m02}{M02. Ноль и разрядный состав}{M02}
\includemodule{module-m03}{M03. Классы многозначных чисел}{M03}
\includemodule{module-m04}{M04. Чтение и запись многозначных чисел}{M04}

\technicalblankpage

\includemodule{module-m05}{M05. Составление чисел из цифр}{M05}
\includemodule{module-m06}{M06. Наибольшее и наименьшее число}{M06}

\technicalblankpage

\includemodule{module-m07}{M07. Упорядоченный перебор и условия между цифрами}{M07}

\technicalblankpage

\end{document}

